(1) 与介质插入方向平行的两块区域相当于并联电容，真空区域面积 $w(L-x)$、介质区域面积 $wx$。故 $C(x)=\varepsilon_0 w(L-x)/d+\kappa\varepsilon_0wx/d=\varepsilon_0w[L+(\kappa-1)x]/d$。在忽略边缘场的近似中其导数 $C'=\varepsilon_0w(\kappa-1)/d>0$，插入介质使电容增大。
(2) 若电荷固定，电场储能 $U_Q=Q^2/[2C(x)]$，机械广义力 $F_x=-dU_Q/dx=Q^2C'(x)/[2C(x)^2]>0$。方向沿更深插入，体现极化介质被吸入强电场区；电容越大，同一固定电荷下的吸引力会逐渐变小。
(3) 电压由电源维持时 $Q(x)=C(x)V$，根据总体系的源—场热力学势 $\mathcal F=U-VQ=-C(x)V^2/2$，得 $F_x=-\mathcal F'(x)=V^2C'/2=\varepsilon_0w(\kappa-1)V^2/(2d)>0$。在微位移中电源向电路输送电量 $dQ=VC'dx$，所作电功为 $dW_{\rm source}=VdQ=V^2C'dx$。
(4) 固定电压时场能增量为 $dU_{\rm field}=(V^2/2)C'dx>0$，介质获得的机械功为 $dW_{\rm mech}=F_xdx=(V^2/2)C'dx$。两者之和恰为电源提供的 $V^2C'dx$，因此能量守恒成立。若误把 $-d(CV^2/2)/dx$ 当作固定电压下的力，会得出排斥的错误符号，是因漏掉电源作功。
